\section{Free Extensions and Free Models/Algebras}
\seclabel{core}

Before delving into the details of \Frexlib{}'s core, we revisit
our frexlet representations using examples for elements in the
fral and the frex for ordinary, commutative, and involutive monoids~(see~\figref{monoid frexlets}).

The elements in the free monoid are lists of the variables appearing
in the term. The elements in the free extension of a
monoid are lists alternating between concrete elements in the given
monoid, and freely adjoined variables. The figure shows an element
in the free extension by 1 variable ($y$) of the multiplicative monoid
of $2\times2$ matrices with natural-number components. The matrix $y$
is unknown, and so its occurrence separates the elements in the list.

Further assuming commutativity equates more terms, resulting
in a representation of the free commutative monoid over \nvar{}
variables as an \nvar-vector of coefficients, representing a linear
polynomial. Freely extending a commutative monoid \A{} by
\nvar{} variables can be represented by a concrete coefficient $c :
\A$ together with an \nvar{}-vector of coefficients, representing a
linear polynomial over $\A$.

If we instead include an involutive operation $x \mapsto
\overline x$ over the monoid, we get lists of variables and alternating lists
whose letters may be tagged as involuted. The figure demonstrates
the free extension of the monoid structure of \String{} concatenation,
with string reversal for the involution.

\subsection{Universal Properties}
We now show the \idris{} realisation of the universal algebraic
concepts from \secref{new-overview}. In the previous \secref{frexlib},
we introduced these definitions:
\begin{itemize}
\item An \IdrisType{Algebra} for a signature, which consists of a carrier
  and an interpretation of operations (\figref{algebras}). It does not
  have to satisfy any equation.
\item A \IdrisType{SetoidAlgebra} is simultaneously a setoid structure
  (an equivalence relation) and an algebra structure over the same
  carrier, whose operations are setoid homomorphisms (\figref{algebras}).
\item A (setoid) \IdrisType{Model} for a presentation is a setoid
  algebra validating the presentation's axioms.
\end{itemize}
In this section, we introduce these definitions:
\begin{itemize}
\item Homomorphisms of setoid algebras.
\item A \IdrisType{ModelOver} a setoid \X{} is a model \A{}  equipped with a setoid
  homomorphism \X{} to \A{}.
\item An \IdrisType{Extension} of a model \A{} by a setoid \X{} is another
  model \A[2]{} together with a setoid homomorphism from \X{} to \A[2]{} and a
  model homomorphism from \A{} to \A[2].
\item Their morphisms.
\end{itemize}
We will then show how we define in \idris{} both the free model
(fral) over a setoid \X{} and the free extension (frex) of \A{} by
\X{} as the initial ones.


\subseclabel{UPs}
\begin{figure}
  \TwoColDIY{.31}{%
    \OpPreservation{}
    }{.65}{%
    \IsAlgHomomorphism{}
    \vspace{1\baselineskip}
  \AlgHomomorphism{}
  }
\caption{Setoid algebra homomorphisms in \Frexlib{}}
\figlabel{homo}
\end{figure}

\Frexlib{} defines homomorphisms of setoid algebras in \figref{homo}, by
requiring the underlying function to be a setoid homomorphism between
the corresponding setoids. The code uses an appropriate
\IdrisFunction{cast} function (see page~\pageref{cast explanation})
that assembles these setoids from the
data in each setoid algebra.
Each $\SigAlgebra$ defines a homomorphic extension operator
\AlgBind{}\AlgBindType{} by structural induction over the
term (i.e.~folding). The following term valuates to \BindExResult{}, e.g.:
\begin{center}
\BindEx{}
\end{center}


Similarly, \figref{algebras and extensions} presents the declarations
for models over a setoid and extension of a model by a setoid. It
expresses the equations in the commuting diagrams through
the extensionality relation on the setoid of functions from
\figref{quotients-and-functions}b in \appref{setoid example} and the power of a
model by a setoid (see \subsecref{powers}). As \idris{}
supports type-directed disambiguation, we overload the record name
\IdrisType{($\~>$)}.

\begin{figure*}
  \TwoColDIY%
      {.43}{%
        \ModelOver{}
        \PreservesEnvDecl{}
        \ModelOverHomo{}
      }
      {.55}{%
        \Extension{}
        \vspace{-.5\baselineskip}
        \ExtensionHomo{}
      }
  \caption{Structure and its preservation for (a) models
    over a setoid, and (b) extensions of a model}
  \figlabel{algebras and extensions}
\end{figure*}



The \emph{free} model over a set (fral) and the \emph{free}
extension (frex) of an algebra by a set is then the initial such
structure: there is a unique structure-preserving map from the free
structure to every structure.
This succinct definition, while
standard, packs much structure. By way of introduction, we will unpack
it for the free commutative monoid over $\Finn$, the finite set with
$\nvar$ elements.

First, we designate a commutative monoid for the
the fral. This structure is the data structure our simplifier will use
to represent the equivalence classes of terms.
In \figref{monoid frexlets}, we mentioned the
carrier consists of origin-intercepting linear polynomials with \Nat{}
coefficients $p = a_1x_1 + \ldots + a_{\nvar}x_{\nvar}$, which we
represent with $\nvar$-tuples of natural numbers and pointwise addition:
\begin{center}
\begin{minipage}{.35\textwidth}
  \ModelCarrier{}
\end{minipage}
\begin{minipage}{.6\textwidth}
  \[
  \begin{aligned}[t]
    0 \definedby{}& 0x_1 +\ldots+0x_n & \quad p \Plus q \definedby{}& (a_1 + b_1)x_1 + \ldots + (a_{\nvar} + b_{\nvar})x_{\nvar} \\
    \definedby{}& \ReplicateZeroPartA \ldots \ReplicateZeroPartB
    & \definedby{}& \OpenBracket a_1 \Plus b_2 \Comma \ldots \Comma a_{\nvar} \Plus b_{\nvar} \CloseBracket
    \\
    ={}& \ReplicateZero & ={}& \ZipIt{}
  \end{aligned}
  \]
\end{minipage}
\end{center}

Denote the resulting $\AppliedFreeCMonoidModelType$ by $\AppliedFreeCMonoidModel$.  For the  $\EnvName$ component, use tabulation to define
$ \FreeCMonoidUnit\,\nvar : \Finn \to \ModelCarrierName$, with $\idrisdata 1$ in
the argument position and $\idrisdata 0$ elsewhere:
%%%% For publication, rename dirac -> kronecker as that is the right
%%%% thing to do after all
\renewcommand\tabulateDirac[1][]{\UseVerb[#1]{tabulateDirac}}
\begin{SaveVerbatim}[commandchars=\\\{\}]{tabulateDirac}
\IdrisFunction{tabulate}\KatlaSpace{}$\KatlaSpace{}\IdrisFunction{kronecker}\KatlaSpace{}\IdrisBound{i}
\end{SaveVerbatim}
\renewcommand\dirac[1][]{\UseVerb[#1]{dirac}}
\begin{SaveVerbatim}[commandchars=\\\{\}]{kronecker}
\IdrisFunction{kronecker}
\end{SaveVerbatim}
\renewcommand\diracij[1][]{\UseVerb[#1]{diracij}}
\begin{SaveVerbatim}[commandchars=\\\{\}]{diracij}
\IdrisFunction{kronecker}\KatlaSpace{}\IdrisBound{i}\KatlaSpace{}\IdrisBound{j}
\end{SaveVerbatim}
\[
\begin{array}[t]{@{}l@{\qquad}l@{}}
\begin{aligned}[t]
\FreeCMonoidUnit\ \nvar\ i \definedby{}& 1x_i \\
\definedby{}& \ReplicateZeroPartA \ldots\KroneckerMiddle\ldots\ReplicateZeroPartB \\
={}& \tabulateDirac
\end{aligned}
&
\begin{array}[t]{@{}l@{}}
\text{where}\\ \diracij \definedby
\begin{cases}
  \diracsi = \diracsj: & \idrisdata{1} \\
  \diracsi \neq \diracsj:  & \idrisdata{0}
\end{cases}
\end{array}
\end{array}
\]
The initiality of this structure follows from the normal form property
--- every origin-intersecting linear polynomial $p$ can be represented
as $p = \sum_{i=1}^n a_i\cdot \FreeCMonoidUnit\ \nvar\ \diracsi$:\\
\[
\begin{array}{@{}l@{}}
  % NB: reverse from the codebase, hopefully that's ok!
  \AppliedCMonoidNormalFormTypePartB
  \\
  \multicolumn{1}{r}{%
    \AppliedCMonoidNormalFormTypePartC\ \AppliedCMonoidNormalFormTypePartA
  }
\end{array}
\]
Since monoid homomorphisms preserve the summation and
multiplication-by-a-natural, the unique structure
preserving map $\hName : \pair{\AppliedFreeCMonoidModel}{\FreeCMonoidUnit\,\nvar} \to \avar$
is this homomorphism:
\[
\hName\ \idrisvar{xs}= \CMonoidUniqueMap{}
\]

This standard argument lies behind many simplifiers, as well as more
advanced techniques like normalisation-by-evaluation. \Frexlib{} takes the same approach, but also explores how to use
general-purpose constructions involving frals and frexes, and bespoke
facts about algebraic structures, to construct new frals and frexes.

To summarise, to implement a fral/frex simplifier, the developer
follows these steps:
\begin{itemize}
\item Design a data-structure for the carrier of the frex/fral's
  algebra, e.g.~for commutative monoids:
  \IdrisType{Vect} \nvar{} \IdrisType{Nat} for the fral
  and \IdrisKeyword{(}\U{} \avar{}, \IdrisType{Vect} \nvar{}
  \IdrisType{Nat}\IdrisKeyword{)} for the frex.
\item Equip it with a setoid algebra structure: pointwise
  operations with propositional equality.
\item Equip it with the appropriate additional structure, e.g.~the
  unit for the fral and the \IdrisFunction{Var}iable function and the
  \IdrisFunction{Embed}ding homomorphism for the frex.
\item Define the function underlying the homomorphism into any other
  algebra over the variable setoid or extension, e.g.~linear
  combination for commutative monoids.
\item Prove that this function is a homomorphism and its uniqueness.
\end{itemize}
These steps realise the simplifier development methodology
we described in \secref{simplification} (page~\pageref{methodology-anchor}).

\subsection{Solver Interface}\subseclabel{solver implementation}
We can now explain the interface \Frexlib{} gives to the \IdrisFunction{solve}
functions. We describe the frex-based interface in detail; the
fral-based interface is similar. We implement the core functionality in the
auxiliary function \IdrisFunction{solveVect} in \figref{solveVect}.
\begin{figure}
  \begin{minipage}{.75\textwidth}
    \solveVectType[numbers=left]{}
  \end{minipage}
  \caption{Core frex-based simplification routine}\figlabel{solveVect}
\end{figure}

The argument \IdrisBound{frex} (line 2) is an implementation of a frex
simplifier for some \IdrisBound{pres}-model \avar{}, extended with
\nvar{} free variables (line 1). We erase the number of variables at
runtime, and so we also erase the type \IdrisData{Fin} \IdrisBound{n}.
We cast an erased type to a setoid instead of an unerased type, i.e.:
\\[.4\baselineskip]
\begin{tabular}{lcr}
  \begin{minipage}[b]{.4\textwidth}
    \irrelevantCast
  \end{minipage}
  &
  instead of
  &
  \begin{minipage}[b]{.4\textwidth}
    \relevantCast
  \end{minipage}
\end{tabular}
\\[.4\baselineskip]
\noindent
The function also takes an \IdrisBound{env}ironment of terms to
substitute for the free variables in the simplification equation (line
2). In this auxiliary function, we present the environment using an
\nvar-ary vector of terms over the algebra's carrier. Next comes the
\IdrisBound{eq}uation we want to discharge (line 3), involving either
concrete values (of type \U{} \avar) and any of the \nvar{} available
variables. Both the \IdrisBound{frex} and the \avar{}lgebra with its
\IdrisBound{env}ironment give rise to extensions in the formal sense,
which we can use to give an environment
for the \IdrisBound{eq}uation in question, namely a setoid
homomorphism from the joint setoid of constants and free variables to
the carrier of the model underlying the extension:
\extEnv{}
where:
\begin{minipage}{.8\textwidth}
  \setoidEither{}
\end{minipage}

We use these environments to interpret the equation, once in the
\IdrisBound{frex} (lines 6--7) and once in the given algebra
(lines~9--10). If the equation holds in all extensions, it will hold
in the \IdrisBound{frex} and in \avar{}, and, moreover, homomorphisms
of extensions will preserve this interpretation.  Interpreting this
equation in the \IdrisBound{frex} may have better decidability
properties over equivalence in \avar{}.

We use \idris{}'s \IdrisKeyword{auto}-implicits mechanism to search
for the equivalence of the \IdrisBound{frex} interpretations. This
mechanism will try to find terms that resolve the implicit argument
\IdrisBound{prf}, using a heuristic informed by unification, that will
also attempt to apply data constructors.

Typically, \idris{}'s judgemental
equality decides
the \IdrisBound{frex} setoid's
equivalence relation, and the number of variables we extend by is known
statically.  This case can happen when the equivalence relation
on the setoid algebra \avar{} is decidable by judgemental
equality. Then, the type of the \IdrisBound{prf} argument (line~5) is
a propositional equality between closed terms. Judgemental
equality decides this relation between the interpretations in the
type of \IdrisBound{prf}.  In \idris{}, the \IdrisKeyword{auto}-search
heuristic tries to use \IdrisData{Refl}, and promotes the required
equation to a judgemental equality constraint.  Even when the setoid relation
is not decidable by judgemental equality,
making \IdrisBound{prf} an \IdrisKeyword{auto}-implicit may provide
more functionality in the future. We may be able to
freely extend algebras whose propositional equality is only partially
decidable by judgemental equality (e.g.~function types in a
type theory with function extensionality), or by a
sophisticated decision procedure (e.g.~multiset equality).

\begin{figure}
  \begin{minipage}{.75\textwidth}
    \Visibility{}
    \metaPi{}
    \metaPI{}
  \end{minipage}
  \caption{Metaprogramming abstractions for curried $\Pi$-types}\figlabel{meta-pi}
\end{figure}

We use metaprogramming abstractions (\figref{meta-pi}) to
simplify the user-facing interface. The \IdrisFunction{PI} combinator
produces an \nvar-ary telescope of
\IdrisData{Visible}/\IdrisData{Hidden}/\IdrisData{Auto} arguments,
packaged as an \nvar-ary vector which it passes this its argument \IdrisBound{b}.
Using this abstraction to reduce
\IdrisFunction{solve} (\figref{solve}) to \IdrisFunction{solveVect} (\figref{solveVect}).
Compare their types: \IdrisFunction{solve} curries the \IdrisBound{env}ironment
$n$-vector into $n$ implicit arguments.
Unification can resolve these arguments to the free
variables in the discharged equation.

\begin{figure}
  \begin{minipage}{.75\textwidth}
    \solveType[numbers=left]{}
  \end{minipage}
  \caption{User-facing frex-based simplification routine}\figlabel{solve}
\end{figure}

\goodbreak
\subsection{Powers}
The commutative monoid structure \AppliedFreeCMonoidModel{} instantiates a general
construction: $\Pres$-models have
powers by setoids. The \emph{power} of an
algebra $\A$ by a set(oid) $\X$ is the terminal
\emph{parameterisation}.
%
\subseclabel{powers}
\begin{wrapfigure}[4]{r}{2.6cm}
  \vspace{-1\baselineskip}
  \includegraphics{diagram-03.mps}
\end{wrapfigure}
%
Parameterisations, shown succinctly on the right, are an $\X$-indexed
collection of algebra homomorphisms $\aEval\ \f : \aModel \to \A$.
%
Requiring $\aEval\ \f$ to be homomorphic implies that operations are
given pointwise. Structure preservation uses the contravariant
action $\idrisfun{pre}\ \UH \hName$ precomposing a homomorphism $\UH
\hName : \aModel \to \bModel$. Universality singles out
the carrier of the power as the function-space \PowerConstruction{}.
For $\X = \Finn$, we can represent it by $\nvar$-tuples from
$\U\,\A$.

%\section{Generic constructions for frexlets}
\subsection{Frex via Coproducts with Fral}
\subseclabel{frex-via-coproducts}
\begin{wrapfigure}[4]{r}{3cm}
  \vspace{-1\baselineskip}
\includegraphics{diagram-06.mps}
\end{wrapfigure}
The fral and the frex are related: the free extension of $\A$ by $\X$ is
the \emph{coproduct} of $\A$ with the free algebra over $\X$.
%
A \emph{cospan} $\avar{}$ between $\A_1$ and $\A_2$ consists of two
homomorphisms with a shared codomain: $\A_1
\xto{\aLft}$\aSink$\xleftarrow{\aRgt} \A_2$.  A \emph{cospan
homomorphism} $h : \avar \to \bvar$ is a homomorphism $\UH : \aSink
\to \bSink$ preserving the cospan as in the diagram on the right.  The
coproducts $\A_1\oplus\A_2$ of two models $\A_1$ and $\A_2$ is then
the initial \emph{cospan}.  All models have coproducts, but these may
be difficult to represent. However, in cases such as commutative
monoids, the coproduct is straightforward to represent: its carrier is
the cartesian product of the components carriers.

The universal property of the frex $\A{}[\X]$ combines
those of the fral $\FreeName\,\Pres\!\X$ and
its coproduct with $\A$. Consider the following two diagrams:
\[
\hfill\includegraphics{diagram-05.mps}
\qquad
\hfill\includegraphics{diagram-04.mps}
\]
The \Var-arrows and the \Lft-arrows correspond through the universal
property of the fral: \Lft{} is the unique homomorphic extension of
the \Var{}-arrows. The \Embed-homomorphisms are exactly
the \Rgt-homomorphism.
This identification lets us construct:
\MyCodeListing{
\FrexByCoprodPartA{}\IdrisKeyword{ :}\FrexByCoprodPartB{}\\
\FrexByCoprodPartC{}
}
For commutative monoids it gives the
commutative monoid of linear polynomials with natural numbers as degree-1
coefficients whose carrier is represented by $\CMonoidFrexCarrier{}$.

\subsection{Fral via a Frex}
\subseclabel{fral via frex}
We use the relationship between the fral and frex, described in \subsecref{frex-via-coproducts}, to derive a fral from a corresponding frex.
%
In particular, the following calculation shows that every free algebra arises as a free extension of $\FreeName\,\Pres\emptyset$. We use $\uplus$ for the disjoint union, i.e., the coproduct of sets:
\begin{align*}
\FreeName\,\Pres\X
&\isomorphic \FreeName\,\Pres(\X\uplus\emptyset) & (\X \isomorphic \X \uplus \emptyset)\\
&\isomorphic (\FreeName\,\Pres\X)\oplus(\FreeName\,\Pres\emptyset) & (\text{$\FreeName\,\Pres$ left adjoint, left adjoints preserve coproducts})\\
&\isomorphic (\FreeName\,\Pres\emptyset){}[\X]
\end{align*}
%
Therefore, we may construct a fral from an initial
algebra and its frex:
\MyCodeListing{\PutListing{.5\textwidth}\ByFrexFull{}}
This generic construction can produce suboptimal
representations.  For example, the initial monoid is easy to
construct: its carrier is the unit type. Freely extending this
initial monoid produces alternating lists, that interleave the unit
value. Taking lists of variables instead leads to a simpler representation but
requires more complicated proofs.  However, generic constructions such as the foregoing allow us to trade efficiency for rapid prototyping.

\subsection{Reusing Frexlets}
\subseclabel{reusing frexlets}
The final example demonstrates reuse of one simplifier when
constructing another. Recall the presentation of involutive monoids
from the end of \secref{frexlib}.
\begin{proposition}[Jacobs]\label{prop:involution}
  The carrier set and monoid operations of the free involutive monoid on \X{}
  are those of the free monoid on the product\/ $\BoolX$.
  Similarly, the frex of an involutive monoid by $\X$ is the
  frex of its underlying monoid by\/ $\BoolX$, extended with an involution operation.
\end{proposition}
We can prove this proposition directly, establishing the involutive
axioms, and have taken this strategy in \Frexlib{}.  However, it is
possible to prove this result without referring to the specific
representation of the monoid fral/frex, resulting in a
simplifier-transformer reusing the code implementing simplifiers for
monoids to implement simplifiers for involutive monoid.  The potential
for this kind of reuse extends beyond monoids.  We can phrase the
result in much greater generality, and give a higher-level proof,
using \citepos{jacobs:involution} axiomatisation of involutions. This
more abstract proof generalises to other notions of involutive
algebras, and we plan to exploit it in the future for generic frexlet
reuse. The more abstract proof goes beyond the scope of this
article, involving more abstract category theoretic notions.

